\scanpage{26}
\deleted{Si $A$ est l'algèbre de von Neumann produit $\unsure{\prod}$ des $A_i$ $(i\in I)$ et d'une algèbre de von Neumann}

Liebe Paula\unsure{!|,}

\deleted{$B$ \inserted{\unsure{sans projecteur}}}

Wieder ein \unsure{Breifr|Breifm},

\deleted{\unsure{minimale}. $T$ \unsure{induit} sur $A_i$ la fonction $\lambda_i T_i$, où $\lambda_i>0$ et $T_i$ est la trace usuelle.}

\deleted{Mais les $A_i$ \unsure{peuvent se plonger dans une}}

\deleted{\unsure{Une}} Une algèbre de von Neumann \deleted{\unsure{quelconque}} $A$ \inserted{\unsure{peut être considérée comme|sans être nécessairement}} \deleted{peut se plonger dans une} algèbre de von Neumann \inserted{\unsure{d'une algèbre sans atomes}} $B$ \unsure{non atomique}, \deleted{\unsure{par exemple}}, et si $A$ est muni d'une trace semi-finie fidèle $T$, on peut choisir un $B$ muni d'une trace semi-finie fidèle qui \unsure{induit} \deleted{\unsure{exactement}} sur $A$ la trace donnée.

Par exemple, on prend l'algèbre des \unsure{applications}, \deleted{\unsure{bornées}} faiblement mesurables et bornées de $[0,1]$ dans $A\subset L(\mathcal H)$, \deleted{\unsure{munie de}} (\unsure{sur}, en réalité sur $\mathcal H\otimes\mathcal H'$) avec la trace
\[
\overline T(f)=\int_0^1 T(f(s))\,ds.
\]
Bien entendu, si $B$ est une algèbre de von Neumann munie d'une trace $T$, et $A$ une sous-algèbre de von Neumann, alors $A\cap\unsure{\mathfrak A_2|\mathfrak B_2}$ \deleted{\unsure{dans}} \unsure{dans l'adhérence} \inserted{\unsure{alors}} l'idéal de définition de la trace sur $A$, $\unsure{\mathfrak H_A|\mu_A}$ est \unsure{la même pour}. \unsure{Considérons} $\unsure{\theta|A}$ \unsure{une rétraction} de $A$ sur $B$, on voit \deleted{que dans les \unsure{formes} positives} \inserted{dans le plongement \unsure{que continuons}, sur les traces} \unsure{finies}, \unsure{tout un théorème encore sur} $A$ est \unsure{analogue}.
\mspage{27}
